% =====================================================================
%  Lecture 05 -- 4G LTE Architecture                      (Persian edition)
% =====================================================================
\documentclass[11pt,a4paper]{article}
\providecommand\ACNroot{../..}
\input{\ACNroot/templates/lecture-fa}

\ACNset{number}{05}
\ACNset{title}{معماری \lr{LTE} نسل چهارم}
\ACNset{shorttitle}{معماری \lr{LTE} نسل چهارم}
\ACNset{titleother}{4G LTE Architecture}
\ACNset{part}{بخش دوم -- دسترسی و هستهٔ نسل چهارم (درس‌های ۵ تا ۱۳)}
\ACNset{time}{۶ تا ۸ ساعت (۳ ساعت مطالعه، ۳ تا ۵ ساعت تمرین)}
\ACNset{level}{متوسط}
\ACNset{prereq}{درس‌های ۱ تا ۴؛ مبانی \lr{OFDM}؛ شبکه‌سازی \lr{IP}}
\ACNset{gen}{4}

\begin{document}
\ACNcover

% ---------------------------------------------------------------------
\section{اهداف یادگیری}
پس از مطالعهٔ این درس باید بتوانید:
\begin{objectives}
  \item محدودیت‌های نسل سوم را که \acro{LTE} را برانگیختند \textbf{توضیح دهید} و الزامات اصلی
        \acro{LTE} را برشمارید.
  \item اصول طراحی \acro{LTE} -- تخت، تمام‌\acro{IP}، کانال‌های مشترک، بدون تحویل نرم -- و دلیل
        مهندسی هر یک را \textbf{بیان کنید}.
  \item سامانهٔ بستهٔ تکامل‌یافته (\lr{E-UTRAN + EPC}) را \textbf{رسم کنید} و همهٔ نقاط مرجع
        (\lr{S1-MME, S1-U, X2, S5/S8, S6a, S11, SGi, Gx, Rx}) را \textbf{نام ببرید}.
  \item بیشینهٔ نرخ فروسوی \acro{LTE} را از شبکهٔ منابع \textbf{استخراج کنید} و یک بودجهٔ تأخیر صفحهٔ
        کاربر \textbf{بسازید}.
  \item نحوهٔ حمل صدا در یک سامانهٔ تمام‌\acro{IP} (\lr{IMS, VoLTE, CS fallback, SRVCC}) را
        \textbf{توصیف کنید}.
  \item جدایی صفحهٔ کنترل و کاربر (\acro{CUPS}) و نقش آن به‌عنوان پلی به سوی \lr{5G} را \textbf{توضیح دهید}.
\end{objectives}

% ---------------------------------------------------------------------
\section{دانش پیش‌نیاز}
باید معماری \acro{UMTS} و جدایی \lr{CS/PS} آن (درس ۳)، تونل \acro{GTP} در \acro{GPRS} (درس ۱) و ایدهٔ
پایهٔ \acro{OFDM} را بدانید: کانال پهن‌باندی که به زیرحامل‌های باریک و متعامد فراوانی تقسیم شده و هر
یک با پیشوند چرخه‌ای در برابر تداخل بین‌نمادی محافظت می‌شود.

\paragraph{واژه‌های کلیدی} واژگان این درس همراه با معادل انگلیسی آن‌ها:
\begin{keyterms}
تکامل بلندمدت & \lr{LTE} & دسترسی رادیویی نسل چهارم \acro{3GPP} (\lr{E-UTRA})، از نسخهٔ ۸. \\
سامانهٔ بستهٔ تکامل‌یافته & \lr{EPS} & سامانهٔ کامل نسل چهارم: \lr{E-UTRAN + EPC}. \\
ایستگاه پایهٔ \acro{LTE} & \lr{eNB (eNodeB)} & ایستگاه پایه همراه با مدیریت منابع رادیویی؛ بدون \acro{RNC}. \\
هستهٔ بستهٔ تکامل‌یافته & \lr{EPC} & هستهٔ تمام‌\acro{IP}: \lr{MME, S-GW, P-GW, HSS, PCRF}. \\
موجودیت مدیریت تحرک & \lr{MME} & گرهٔ صفحهٔ کنترل: \acro{NAS}، تحرک، امنیت، کنترل حامل. \\
دروازهٔ خدمت‌دهنده / دروازهٔ \acro{PDN} & \lr{S-GW / P-GW} & لنگرهای صفحهٔ کاربر: لنگر تحرک / لنگر \acro{IP}. \\
حامل \acro{EPS} & \lr{EPS bearer} & لولهٔ کیفیت خدمات سرتاسری میان \acro{UE} و \lr{P-GW}. \\
دسترسی چندگانهٔ متعامد / تک‌حامل & \lr{OFDMA / SC-FDMA} & دسترسی چندگانهٔ فروسو / فراسوی \acro{LTE}. \\
صدا روی \acro{LTE} & \lr{VoLTE} & صدای مبتنی بر \acro{IMS} روی حامل \acro{EPS} با \lr{QCI 1}. \\
جدایی صفحهٔ کنترل و کاربر & \lr{CUPS} & تقسیم \lr{S/P-GW} به بخش‌های \lr{C} و \lr{U} در نسخهٔ ۱۴ (\acro{PFCP}). \\
\end{keyterms}

% ---------------------------------------------------------------------
\section{مقدمه و انگیزه}
تا سال ۲۰۰۴ \acro{HSPA} پهن‌باند موبایل را ارائه می‌کرد، اما پایه‌هایش کهنه می‌شد. حامل
\qty{5}{\mega\hertz} \acro{WCDMA} نمی‌توانست بدون همسان‌سازهای پیچیده به ۱۰ تا ۲۰ مگاهرتز موردنیاز
برای نرخ‌های بالاتر برسد؛ سلسله‌مراتب \lr{UE--NodeB--RNC--SGSN--GGSN} تأخیر و هزینه می‌افزود؛ دو
هستهٔ موازی (\acro{CS} و \acro{PS}) بهره‌برداری را دو برابر می‌کرد. در همین حال فناوری‌های رقیب
(\acro{WiMAX}) پهن‌باند تمام‌\acro{IP} مبتنی بر \acro{OFDM} را وعده می‌دادند.

مطالعهٔ تکامل بلندمدت \acro{3GPP} اهداف بلندپروازانه‌ای تعیین کرد \cite{tr25913}: بیشینهٔ نرخ
\qty{100}{\mega\bit\per\second} فروسو و \qty{50}{\mega\bit\per\second} فراسو در \qty{20}{\mega\hertz}؛
تأخیر صفحهٔ کاربر کمتر از ۵ میلی‌ثانیه؛ گذار از بیکار به فعال کمتر از ۱۰۰ میلی‌ثانیه؛ بازدهی طیفی ۲
تا ۴ برابر \acro{HSPA} نسخهٔ ۶؛ پهنای باند مقیاس‌پذیر از ۱٫۴ تا ۲۰ مگاهرتز؛ و پشتیبانی از سرعت تا ۳۵۰
کیلومتر بر ساعت. \emph{تکامل معماری سامانه} (\acro{SAE}) یک هستهٔ متناظر، یعنی \acro{EPC}، را تعریف
کرد. این دو با هم سامانهٔ بستهٔ تکامل‌یافته (\acro{EPS}) را تشکیل می‌دهند که نخستین بار در نسخهٔ ۸
(۲۰۰۸) مشخص شد \cite{ts36300,ts23401}.

% ---------------------------------------------------------------------
\section{مفاهیم بنیادی}

\subsection{اصول طراحی و دلایل آن‌ها}
\begin{itemize}
  \item \textbf{شبکهٔ دسترسی تخت.} \acro{RNC} حذف می‌شود؛ کارکردهایش به \acro{eNB} (کنترل منابع
        رادیویی، زمان‌بندی، رمزنگاری، تصمیم تحویل) یا به هسته منتقل می‌شوند. گره‌های کمتر یعنی تأخیر
        و هزینهٔ کمتر. بهای آن این است که هیچ گره‌ای نمی‌تواند سیگنال چند سلول را ترکیب کند -- بنابراین
        تحویل نرم وجود ندارد (درس ۴).
  \item \textbf{هستهٔ تمام‌بسته‌ای.} هیچ حوزهٔ سوئیچ مداری وجود ندارد. صدا باید به‌صورت بسته‌های
        \acro{IP} (\acro{VoLTE} روی \acro{IMS}) حمل شود یا موقتاً به نسل دوم/سوم سپرده شود (\lr{CS fallback}).
  \item \textbf{فقط کانال‌های مشترک.} همهٔ دادهٔ کاربر از کانال‌های مشترکی استفاده می‌کند که \acro{eNB}
        هر ۱ میلی‌ثانیه زمان‌بندی می‌کند؛ هیچ کانال ترافیک اختصاصی وجود ندارد.
  \item \textbf{\acro{OFDMA} در فروسو و \acro{SC-FDMA} در فراسو.} \acro{OFDM} به‌خوبی با پهنای باند
        مقیاس می‌شود و زمان‌بندی گزینشی در فرکانس و \lr{MIMO} را ساده می‌کند. در فراسو، \acro{FDMA}
        تک‌حامل نسبت توان قله به میانگین را برای تقویت‌کنندهٔ گوشی پایین نگه می‌دارد (درس ۷).
  \item \textbf{جدایی صفحهٔ کنترل و کاربر} در هسته: \acro{MME} هرگز به بسته‌های کاربر دست نمی‌زند؛
        دروازه‌ها هرگز سیگنالینگ \acro{NAS} را پایان نمی‌دهند.
  \item \textbf{شبکه‌های خودسازمان‌ده (\acro{SON}).} روابط خودکار همسایگی، استحکام تحرک و توازن بار،
        برنامه‌ریزی دستی را در شبکه‌ای از گره‌های فراوان، کوچک و تخت کاهش می‌دهند.
\end{itemize}

\subsection{حامل \lr{EPS}}
\begin{definitionbox}{حامل \lr{EPS}}
حامل \acro{EPS} لوله‌ای منطقی میان \acro{UE} و \lr{P-GW} با یک رفتار کیفیت خدمات واحد است. از به هم
پیوستن یک \emph{حامل رادیویی} (\lr{UE--eNB})، یک \emph{حامل \lr{S1}} (\lr{eNB--S-GW}، یک تونل
\lr{GTP-U}) و یک \emph{حامل \lr{S5/S8}} (\lr{S-GW--P-GW}، تونل \lr{GTP-U} دیگر) تشکیل می‌شود. هر
\acro{UE} هنگام اتصال یک \emph{حامل پیش‌فرض} می‌گیرد؛ \emph{حامل‌های اختصاصی} برای ترافیکی که
کیفیت خدمات متفاوتی لازم دارد، مانند \acro{VoLTE} (\lr{QCI 1})، افزوده می‌شوند.
\end{definitionbox}
برخلاف \acro{GPRS} که اتصال‌پذیری پس از اتصال به فعال‌سازی جداگانهٔ زمینهٔ \acro{PDP} نیاز داشت،
\acro{LTE} حامل پیش‌فرض را \emph{در حین} اتصال برقرار می‌کند: دستگاه \acro{LTE} همیشه روشن و دارای
نشانی \acro{IP} است. حامل‌ها و کلاس‌های کیفیت خدمات در درس‌های ۹ و ۲۲ گسترش می‌یابند.

% ---------------------------------------------------------------------
\section{توضیحات فنی تفصیلی}

\subsection{عناصر شبکه}
\paragraph{\acro{UE}} تجهیزات موبایل به‌اضافهٔ \acro{UICC} (با برنامهٔ \acro{USIM}). \acro{UE}های
\acro{LTE} بر اساس \emph{رده}هایی دسته‌بندی می‌شوند که بیشینهٔ نرخ، لایه‌های \lr{MIMO} و اندازهٔ
بافر را مشخص می‌کنند.

\paragraph{\acro{eNB}} تنها گرهٔ \acro{E-UTRAN} \cite{ts36300}. کل پشتهٔ رادیویی (\lr{PHY, MAC, RLC,
PDCP, RRC}) را اجرا می‌کند، فراسو و فروسو را هر میلی‌ثانیه زمان‌بندی می‌کند و کنترل پذیرش، تصمیم
تحویل، فشرده‌سازی سرآیند و رمزنگاری دادهٔ کاربر را انجام می‌دهد، دادهٔ کاربر را به \lr{S-GW} مسیریابی
می‌کند و هنگام اتصال برای \acro{UE} یک \acro{MME} انتخاب می‌کند. \acro{eNB}ها از طریق \acro{X2} به
یکدیگر و از طریق \lr{S1} به هسته وصل می‌شوند.

\paragraph{\acro{MME}} مغز صفحهٔ کنترل \acro{EPC}: سیگنالینگ \acro{NAS} و امنیت آن را پایان می‌دهد،
هویت \acro{UE} را با بردارهای دریافتی از \acro{HSS} احراز می‌کند، نواحی ردیابی و فراخوانی را مدیریت
می‌کند، \lr{S-GW} و \lr{P-GW} را انتخاب می‌کند و حامل‌ها را از طریق \lr{S-GW} برقرار، اصلاح و آزاد
می‌کند. موضوع درس ۱۲ است.

\paragraph{\lr{S-GW}} \lr{S1-U} را پایان می‌دهد و \emph{لنگر تحرک محلی} است: در تحویل میان \acro{eNB}ها
تنها تونل \lr{S1-U} تغییر می‌کند. دادهٔ فروسو را برای \acro{UE}های بیکار بافر می‌کند و از طریق \acro{MME}
فراخوانی را آغاز می‌کند، و تحرک به نسل دوم و سوم را لنگر می‌کند.

\paragraph{\lr{P-GW}} \emph{لنگر \acro{IP}}: نشانی \acro{IP} \acro{UE} را تخصیص می‌دهد، نقطهٔ اتصال به یک
شبکهٔ دادهٔ بسته‌ای (که با یک \acro{APN} شناسایی می‌شود) است، سیاست و صورت‌حساب را اعمال می‌کند
(\acro{PCEF})، بازرسی عمیق بسته و شنود قانونی انجام می‌دهد. یک \acro{UE} می‌تواند برای \acro{APN}های
مختلف (مثلاً \lr{\texttt{internet}} و \lr{\texttt{ims}}) به چند \lr{P-GW} وصل باشد. درس ۱۳ هر دو دروازه
را بررسی می‌کند.

\paragraph{\acro{HSS}} پایگاه دادهٔ مشترکان، تکامل‌یافتهٔ \acro{HLR}: نمایهٔ اشتراک، کارکردهای
\acro{AuC} و \acro{MME} فعلی. درس ۱۱.

\paragraph{\acro{PCRF}} تابع قواعد سیاست و صورت‌حساب \cite{ts23203}: قواعد کیفیت خدمات و صورت‌حساب را
برای هر جریان دادهٔ خدمت تعیین و از طریق \acro{Gx} به \lr{P-GW} ارسال می‌کند؛ توابع کاربردی مانند
\lr{P-CSCF} در \acro{IMS} از طریق \acro{Rx} منابع درخواست می‌کنند.

\subsection{نقاط مرجع}
\begin{table}[htbp]
\centering\small
\caption{نقاط مرجع اصلی \acro{EPS} \cite{ts23401}.}\label{tab:ref}
\rowcolors{2}{acnmist}{white}
\begin{tabularx}{\linewidth}{@{}l l l Y@{}}
\toprule
\rowcolor{white}نقطهٔ مرجع & میان & پروتکل & کاربرد \\
\midrule
\lr{LTE-Uu}  & \lr{UE -- eNB}     & \lr{RRC, PDCP/RLC/MAC/PHY} & واسط رادیویی \\
\lr{S1-MME}  & \lr{eNB -- MME}    & \lr{S1AP over SCTP}        & زمینهٔ \acro{UE}، حامل‌ها، فراخوانی، تحویل \\
\lr{S1-U}    & \lr{eNB -- S-GW}   & \lr{GTP-U over UDP/IP}     & دادهٔ کاربر به ازای هر حامل \\
\lr{X2}      & \lr{eNB -- eNB}    & \lr{X2AP / GTP-U}          & تحویل، ارسال داده، هماهنگی تداخل \\
\lr{S11}     & \lr{MME -- S-GW}   & \lr{GTPv2-C}               & مدیریت نشست و حامل \\
\lr{S5/S8}   & \lr{S-GW -- P-GW}  & \lr{GTPv2-C + GTP-U}       & \lr{S5} درون یک شبکه، \lr{S8} میان شبکهٔ خانگی و میزبان \\
\lr{S6a}     & \lr{MME -- HSS}    & \lr{Diameter}              & بردارهای احراز هویت، به‌روزرسانی مکان، اشتراک \\
\lr{SGi}     & \lr{P-GW -- PDN}   & \lr{IP}                    & به اینترنت، \acro{IMS} یا شبکه‌های سازمانی \\
\lr{Gx}      & \lr{PCRF -- P-GW}  & \lr{Diameter}              & قواعد سیاست و صورت‌حساب \\
\lr{Rx}      & \lr{AF -- PCRF}    & \lr{Diameter}              & درخواست کیفیت خدمات از سوی کاربرد (مثلاً \acro{VoLTE}) \\
\lr{Gy/Gz}   & \lr{P-GW -- OCS/OFCS} & \lr{Diameter / file}    & صورت‌حساب برخط / برون‌خط \\
\bottomrule
\end{tabularx}
\end{table}

دو الگو برجسته است. نخست، هر واسط کنترلی هسته یا \lr{GTPv2-C} است (میان گره‌هایی که تونل‌ها را
مدیریت می‌کنند) یا \lr{Diameter} (به سوی پایگاه‌های داده و سیاست) -- \acro{SS7} و \acro{MAP} از
\acro{EPC} ناپدید شده‌اند. دوم، \lr{S8} گونهٔ \emph{رومینگ} \lr{S5} است: با \emph{مسیریابی خانگی}،
\lr{S-GW} میزبان ترافیک را به \lr{P-GW} خانگی تونل می‌کند تا اپراتور خانگی کنترل نشانی‌دهی \acro{IP} و
صورت‌حساب را حفظ کند.

\subsection{صدا در یک سامانهٔ تمام‌\lr{IP}}
\acro{LTE} بدون حوزهٔ \acro{CS} سه راه‌حل صوتی ارائه می‌دهد:
\begin{itemize}
  \item \textbf{\acro{VoLTE}} -- صدا به‌عنوان یک خدمت \acro{IMS} \cite{ts23228}. \acro{UE} روی یک حامل
        پیش‌فرض به \acro{APN} \lr{\texttt{ims}} در \acro{IMS} ثبت‌نام می‌کند (\lr{QCI 5} برای سیگنالینگ
        \acro{SIP}). هنگام برقراری تماس، \lr{P-CSCF} از طریق \acro{Rx} از \acro{PCRF} منابع درخواست
        می‌کند؛ \acro{PCRF} از طریق \acro{Gx} به \lr{P-GW} دستور می‌دهد یک حامل اختصاصی \acro{GBR} با
        \lr{QCI 1} برای بسته‌های گفتار بسازد (درس ۱۰).
  \item \textbf{بازگشت به \acro{CS} (\acro{CSFB})} -- \acro{UE} از طریق \acro{MME} (واسط \lr{SGs}) در حوزهٔ
        \acro{CS} نسل دوم/سوم ثبت می‌شود؛ برای تماس به نسل دوم/سوم هدایت می‌شود \cite{ts23272}. ساده
        است، اما برقراری تماس کندتر است و در طول تماس سرعت داده کاهش می‌یابد.
  \item \textbf{\acro{SRVCC}} -- پیوستگی تماس صوتی تک‌رادیویی: تماس \acro{VoLTE} در جریان، هنگامی که
        \acro{UE} پوشش \acro{LTE} را ترک می‌کند، به حوزهٔ \acro{CS} تحویل می‌شود \cite{ts23216}.
\end{itemize}

\subsection{جدایی صفحهٔ کنترل و کاربر (\lr{CUPS})}
در نسخهٔ ۸، \lr{S-GW} و \lr{P-GW} هر یک هم منطق کنترل (\lr{GTP-C}، سیاست، تخصیص \acro{IP}) و هم ارسال
بسته را در بر داشتند. نسخهٔ ۱۴ آن‌ها را تقسیم کرد \cite{ts23214}: \lr{SGW-C} و \lr{PGW-C} منطق کنترل را
نگه می‌دارند و \lr{SGW-U} و \lr{PGW-U} را از طریق واسط \lr{Sx} و با پروتکل کنترل ارسال بسته (\acro{PFCP})
\cite{ts29244} برنامه‌ریزی می‌کنند. صفحه‌های کاربر را می‌توان مستقل مقیاس داد و برای تأخیر کم نزدیک
رادیو قرار داد (رایانش لبه)، در حالی که کنترل متمرکز می‌ماند. \acro{CUPS} نیای مستقیم تقسیم
\lr{SMF/UPF} در \lr{5G} است (درس‌های ۱۳ و ۱۹).

% ---------------------------------------------------------------------
\section{مدل‌های ریاضی}

\subsection{بیشینهٔ نرخ فروسو از شبکهٔ منابع}\label{sec:peak}
\paragraph{مسئله} «۳۰۰ مگابیت بر ثانیهٔ» یک دستگاه \acro{LTE} ردهٔ ۵ از کجا می‌آید؟
\paragraph{مدل سامانه} پهنای باند \qty{20}{\mega\hertz}، پیشوند چرخه‌ای عادی، \acro{FDD}. شبکهٔ منابع
(درس ۷) دارای $N_{\mathrm{RB}} = 100$ بلوک منبع با ۱۲ زیرحامل با $\Delta f = 15$\,\lr{kHz} و ۱۴ نماد
\acro{OFDM} در هر زیرقاب ۱ میلی‌ثانیه‌ای است، یعنی $12\times14 = 168$ عنصر منبع (\acro{RE}) به ازای هر
\acro{RB} در هر زیرقاب. فرض‌ها: یک نماد \acro{OFDM} برای ناحیهٔ کنترل، چهار درگاه سیگنال مرجع ویژهٔ
سلول، \lr{64QAM} ($Q_m = 6$)، $\nu = 4$ لایهٔ \lr{MIMO}؛ کانال‌های همگام‌سازی و پخش را که تنها چند
زیرقاب را اشغال می‌کنند نادیده می‌گیریم.
\paragraph{استخراج} چهار درگاه \acro{CRS} در هر جفت \acro{RB} از ۲۴ عنصر منبع استفاده می‌کنند که چهار عدد
از آن‌ها در نماد کنترل قرار دارند؛ بنابراین عناصر منبع داده به ازای هر \acro{RB} در هر زیرقاب:
\begin{equation}
  N_{\mathrm{RE}}^{\mathrm{data}} = 168 - \underbrace{12}_{\text{control}} - \underbrace{20}_{\text{CRS}} = 136 .
\end{equation}
بیشینهٔ نرخ برابر است با
\begin{equation}\label{eq:peak}
  R_{\mathrm{peak}} = \nu\, Q_m\, r_\mathrm{c}\,\frac{N_{\mathrm{RB}}\,N_{\mathrm{RE}}^{\mathrm{data}}}{T_{\mathrm{sf}}}
  = 4\times6\times r_\mathrm{c}\times\frac{100\times136}{\qty{1}{\milli\second}}
  = 326.4\,r_\mathrm{c}\ \mathrm{Mbit/s},
\end{equation}
که $r_\mathrm{c}$ نرخ کد کانال و $T_{\mathrm{sf}}$ مدت زیرقاب است. با بزرگ‌ترین بلوک انتقال استانداردشده
($r_\mathrm{c}\approx0.92$) حاصل حدود $\qty{300}{\mega\bit\per\second}$ است که با
\lr{299.6\,Mbit/s} ردهٔ ۵ \acro{UE} سازگار است.
\paragraph{تفسیر} سیگنال‌های مرجع و ناحیهٔ کنترل حدود ۱۹٪ عناصر منبع را مصرف می‌کنند؛ پیشوند چرخه‌ای
(۶٫۷٪ زمان، \cref{exr:cp}) از پیش در زیرقاب ۱۴ نمادی گنجانده شده است. بیشینهٔ نرخ به چهار آنتن
گیرنده، \acro{SINR} عالی و سلول خالی نیاز دارد -- کرانی بالایی برای واسط هوایی است، نه تجربهٔ کاربر.

\subsection{بودجهٔ تأخیر صفحهٔ کاربر}\label{sec:latency}
\paragraph{مسئله} آیا \acro{LTE} می‌تواند هدف تأخیر یک‌طرفهٔ ۵ میلی‌ثانیه‌ای صفحهٔ کاربر را برآورده کند؟
\paragraph{مدل} تأخیر یک‌طرفه از رسیدن یک بسته به \acro{PDCP} در \acro{eNB} تا تحویل آن به لایهٔ \acro{IP}
در \acro{UE}، در سلول بدون بار، به این صورت تجزیه می‌شود:
\begin{equation}\label{eq:latency}
  D_{\mathrm{UP}} = \underbrace{\tfrac{1}{2}T_{\mathrm{TTI}}}_{\text{alignment}}
  + T_{\mathrm{TTI}} + D_{\mathrm{eNB}} + D_{\mathrm{UE}} + p_{\mathrm{re}}\,T_{\mathrm{HARQ}}
  + D_{\mathrm{tr}},
\end{equation}
که جملهٔ هم‌ترازی میانگین انتظار تا مرز \acro{TTI} بعدی، $D_{\mathrm{eNB}}$ و $D_{\mathrm{UE}}$ تأخیرهای
پردازش، $p_{\mathrm{re}}$ احتمال یک ارسال مجدد \acro{HARQ}، $T_{\mathrm{HARQ}} = \qty{8}{\milli\second}$
زمان رفت‌وبرگشت \acro{HARQ} در \acro{FDD} و $D_{\mathrm{tr}}$ تأخیر باقی‌ماندهٔ انتقال است.
\paragraph{مثال عددی} با $T_{\mathrm{TTI}} = 1$\,\lr{ms}، $D_{\mathrm{eNB}} = D_{\mathrm{UE}} = 1$\,\lr{ms}،
$p_{\mathrm{re}} = 0.1$ و $D_{\mathrm{tr}} = 0.5$\,\lr{ms} (مقادیر نمایشی، \cref{fig:latency}):
$D_{\mathrm{UP}} = 0.5 + 1 + 1 + 1 + 0.8 + 0.5 = \qty{4.8}{\milli\second}$.
\paragraph{تفسیر و محدودیت‌ها} \acro{TTI} یک‌میلی‌ثانیه‌ای و حلقهٔ \acro{HARQ} هشت‌میلی‌ثانیه‌ای غالب‌اند:
این‌ها اهرم‌هایی هستند که \lr{5G} به کار می‌گیرد (شیارهای کوتاه‌تر، ریزشیارها، \acro{HARQ} سریع‌تر، درس
۱۴). مدل صف (بار صفر)، تأخیر درخواست زمان‌بندی در فراسو (\acro{UE} بدون مجوز باید ابتدا درخواست کند که
چند میلی‌ثانیه می‌افزاید) و تأخیر هسته و اینترنت را نادیده می‌گیرد. به همین دلیل زمان‌های \lr{ping}
سرتاسری اندازه‌گیری‌شده در شبکه‌های \acro{LTE} معمولاً ده‌ها میلی‌ثانیه است، هرچند هدف رادیویی ۵
میلی‌ثانیه است.

\begin{acnfigure}
  \acnfig{lecture-05/latency-budget}
  \caption{بودجهٔ نمایشی تأخیر یک‌طرفهٔ صفحهٔ کاربر فروسو بر اساس \cref{eq:latency} (\lr{4.8\,ms}).}
  \label{fig:latency}
\end{acnfigure}

\subsection{سربار بسته‌بندی}\label{sec:overhead}
هر بستهٔ کاربر که از \lr{S1-U} یا \lr{S5} عبور می‌کند در \lr{GTP-U/UDP/IPv4} حمل می‌شود که
$h = 8 + 8 + 20 = 36$ بایت می‌افزاید. برای محمولهٔ $L$ بایتی، بازده انتقال برابر است با
\begin{equation}\label{eq:eff}
  \eta_{\mathrm{tr}} = \frac{L}{L + h}.
\end{equation}
یک بستهٔ وب ۱۴۰۰ بایتی $\eta_{\mathrm{tr}} = 97.5\%$ دارد؛ یک بستهٔ \acro{VoLTE} شصت‌بایتی (قاب گفتار
به‌علاوهٔ سرآیندهای \lr{RTP/UDP/IP}؛ فشرده‌سازی سرآیند تنها روی پیوند رادیویی عمل می‌کند) تنها ۶۲٫۵٪.
ابعادبندی شبکهٔ انتقال برای صدا باید این را در نظر بگیرد (\cref{exr:volte})؛ انتقال \lr{IPv6} ۲۰ بایت دیگر
می‌افزاید.

% ---------------------------------------------------------------------
\section{معماری شبکه و نمودارها}
\Cref{fig:arch} سامانهٔ \acro{EPS} را نشان می‌دهد. در مقایسه با \acro{UMTS} (درس ۳، شکل ۳) شبکهٔ دسترسی
یک نوع گره دارد، هسته حوزهٔ \acro{CS} ندارد و صفحهٔ کاربر تنها از دو گرهٔ هسته (\lr{S-GW} و \lr{P-GW})
عبور می‌کند که هر دو می‌توانند هم‌مکان باشند. صفحه‌های کنترل (خط‌چین) و کاربر (فیروزه‌ای توپر) همه‌جا
جز رادیو از واسط‌های متفاوت استفاده می‌کنند. \Cref{fig:stack} پشتهٔ صفحهٔ کاربر را نشان می‌دهد:
پروتکل‌های رادیویی در \acro{eNB} پایان می‌یابند و از آنجا تونل‌های \lr{GTP-U} بسته‌های \acro{IP} کاربر
را حمل می‌کنند.

\begin{acnfigure}
  \resizebox{\linewidth}{!}{\acnfig{lecture-05/lte-architecture}}
  \caption{سامانهٔ بستهٔ تکامل‌یافته. آبی خط‌چین: صفحهٔ کنترل؛ فیروزه‌ای توپر: صفحهٔ کاربر.}
  \label{fig:arch}
\end{acnfigure}

\begin{acnfigure}
  \resizebox{\linewidth}{!}{\acnfig{lecture-05/up-stack}}
  \caption{پشتهٔ پروتکل صفحهٔ کاربر \acro{LTE}. بستهٔ \acro{IP} \acro{UE} روی رادیو با
  \lr{PDCP/RLC/MAC} و روی \lr{S1-U} و \lr{S5/S8} با تونل‌های \lr{GTP-U} حمل می‌شود؛ \lr{P-GW} لنگر \acro{IP} است.}
  \label{fig:stack}
\end{acnfigure}

\begin{acnfigure}
  \acnfig{lecture-05/cups}
  \caption{جدایی صفحهٔ کنترل و کاربر در دروازه‌های \acro{EPC} (نسخهٔ ۱۴).}
  \label{fig:cups}
\end{acnfigure}

% ---------------------------------------------------------------------
\section{رویه‌های گام‌به‌گام}
\begin{procedurebox}{اتصال اولیه در \lr{LTE} (نمای کلی)}\label{proc:attach}
\begin{steps}
  \item \acro{UE} دسترسی تصادفی انجام می‌دهد و با \acro{eNB} یک اتصال \acro{RRC} برقرار می‌کند (درس‌های ۶ و ۸).
  \item پیام \acro{NAS} \msg{Attach Request} را همراه با \msg{PDN Connectivity Request} می‌فرستد؛
        \acro{eNB} یک \acro{MME} انتخاب و آن را در \lr{S1AP} \msg{Initial UE Message} ارسال می‌کند.
  \item \acro{MME} هویت \acro{UE} را با \lr{EPS-AKA} (بردارها از \acro{HSS} روی \lr{S6a}) احراز می‌کند،
        امنیت \acro{NAS} را فعال می‌کند و خود را در \acro{HSS} ثبت می‌کند (\msg{Update Location}) و
        اشتراک را دریافت می‌کند.
  \item \acro{MME} یک \lr{S-GW} و یک \lr{P-GW} انتخاب و \lr{GTPv2-C} \msg{Create Session Request} را
        می‌فرستد؛ \lr{P-GW} یک نشانی \acro{IP} تخصیص می‌دهد، سیاست را از \acro{PCRF} می‌گیرد و حامل
        پیش‌فرض را می‌سازد.
  \item \acro{MME} پیام \lr{S1AP} \msg{Initial Context Setup Request} را همراه با \msg{Attach Accept}
        می‌فرستد؛ \acro{eNB} امنیت لایهٔ دسترسی را فعال و حامل رادیویی را با
        \msg{RRC Connection Reconfiguration} پیکربندی می‌کند.
  \item \acro{UE} با \msg{Attach Complete} پاسخ می‌دهد؛ \acro{MME} پیام \msg{Modify Bearer Request} را
        می‌فرستد تا \lr{S-GW} نقطهٔ پایانی تونل \acro{eNB} را بداند. دادهٔ فراسو و فروسو می‌تواند جریان یابد.
\end{steps}
\end{procedurebox}
جریان کامل پیام‌ها، با هر عنصر اطلاعاتی مهم، در درس‌های ۹، ۱۱ و ۱۲ ساخته می‌شود. از هم‌اکنون توجه
کنید که یک رویه ثبت‌نام، امنیت \emph{و} نشانی \acro{IP} را فراهم می‌کند -- سه رویهٔ جدای \acro{GPRS}
در یک رویه ادغام شده‌اند.

% ---------------------------------------------------------------------
\section{مثال‌های حل‌شده}

\begin{examplebox}{بیشینهٔ نرخ یک سلول \lr{10\,MHz}}\label{ex:peak10}
\textbf{مسئله.} یک حامل \acro{LTE} ده‌مگاهرتزی (۵۰ \acro{RB}) از دو درگاه \acro{CRS}، دو نماد کنترل و
\lr{MIMO} $2\times2$ با \lr{64QAM} استفاده می‌کند. بیشینهٔ نرخ را برای $r_\mathrm{c} = 0.9$ تخمین بزنید.

\textbf{حل.} دو درگاه \acro{CRS} در هر جفت \acro{RB} از ۱۶ عنصر منبع استفاده می‌کنند که ۴ عدد از آن‌ها در
نخستین نماد (کنترلی) قرار دارند؛ با دو نماد کنترل، $N_{\mathrm{RE}}^{\mathrm{data}} = 168 - 24 - 12 = 132$.
$R = 2\times6\times0.9\times50\times132/10^{-3} = \qty{71.3}{\mega\bit\per\second}$.

\textbf{تفسیر.} نصف پهنای باند و نصف لایه‌ها تقریباً یک‌چهارم رقم ردهٔ ۵ را می‌دهد؛ سربار ناحیهٔ کنترل و
\acro{RS} مشابه است (حدود ۲۱٪).
\end{examplebox}

\begin{examplebox}{گره‌های صفحهٔ کاربر}\label{ex:hops}
\textbf{مسئله.} گره‌های شبکه‌ای را که یک بستهٔ کاربر را میان رادیو و اینترنت در \lr{UMTS R99} و \acro{LTE}
پردازش می‌کنند بشمارید و اگر هر گره ۰٫۵ میلی‌ثانیه پردازش و صف بیفزاید، صرفه‌جویی را تخمین بزنید.

\textbf{حل.} \acro{UMTS}: \lr{NodeB, RNC, SGSN, GGSN} (۴ گره). \acro{LTE}: \lr{eNB, S-GW, P-GW} (۳ گره؛ اگر
\lr{S-GW} و \lr{P-GW} ترکیب شوند ۲ گره). صرفه‌جویی: ۰٫۵ تا ۱ میلی‌ثانیه یک‌طرفه و ۱ تا ۲ میلی‌ثانیه رفت‌وبرگشت.

\textbf{تفسیر.} تعداد گره‌ها اهمیت کمتری از \acro{TTI}های طولانی و زمان‌بندی کند \acro{RNC} دارد که
\acro{eNB} آن را با زمان‌بندی یک‌میلی‌ثانیه‌ای جایگزین می‌کند -- دلیل اصلی اینکه تأخیر \acro{LTE} کسری از
تأخیر \lr{R99} است.
\end{examplebox}

\begin{examplebox}{ظرفیت \lr{VoLTE} یک سلول}\label{ex:volte}
\textbf{مسئله.} هر تماس \acro{VoLTE} به‌طور میانگین در هر جهت در هر ۲۰ میلی‌ثانیه به ۲ \acro{RB} نیاز دارد
(یک قاب گفتار در هر ۲۰ میلی‌ثانیه، \acro{MCS} مقاوم). اگر ۷۰٪ \acro{RB}ها برای صدا قابل استفاده باشد، یک
سلول ده‌مگاهرتزی چند تماس را حمل می‌کند؟

\textbf{حل.} هر تماس: ۲ \acro{RB} در ۲۰ زیرقاب $=$ ۰٫۱ \acro{RB} در هر زیرقاب. در دسترس: $0.7\times50 = 35$
\acro{RB} در هر زیرقاب. تماس‌ها $= 35/0.1 = 350$ در هر جهت (در عمل کانال کنترل \acro{PDCCH} زودتر محدودیت
ایجاد می‌کند و به همین دلیل برای \acro{VoLTE} از \emph{زمان‌بندی نیمه‌پایدار} استفاده می‌شود).

\textbf{تفسیر.} مقایسه کنید با حدود ۴۵ کاربر صوتی \acro{WCDMA} در هر ۵ مگاهرتز (درس ۳): صدای بسته‌ای با
فشرده‌سازی سرآیند و زمان‌بندی، از نظر طیفی بسیار کاراتر از صدای مداری است.
\end{examplebox}

% ---------------------------------------------------------------------
\section{مقایسهٔ فناوری‌ها}
\begin{table}[htbp]
\centering\small
\caption{نسل سوم (\lr{UMTS/HSPA}) در برابر نسل چهارم (\lr{LTE/EPC}).}\label{tab:3g4g}
\rowcolors{2}{acnmist}{white}
\begin{tabularx}{\linewidth}{@{}l Y Y@{}}
\toprule
\rowcolor{white} & \lr{UMTS / HSPA} & \lr{LTE / EPC} \\
\midrule
دسترسی رادیویی       & \acro{WCDMA}، \lr{5\,MHz}            & \acro{OFDMA} فروسو / \acro{SC-FDMA} فراسو، \lr{1.4--20\,MHz} \\
گره‌های شبکهٔ دسترسی & \lr{NodeB + RNC}                     & فقط \acro{eNB} (تخت) \\
تحویل                & نرم + سخت                            & سخت (با ارسال داده روی \acro{X2}) \\
حوزه‌های هسته        & \lr{CS + PS}                         & فقط \acro{PS} \\
صدا                  & \acro{CS}                            & \acro{VoLTE} (\acro{IMS})؛ \lr{CSFB/SRVCC} برای هم‌کنشی \\
سیگنالینگ هسته       & \lr{MAP (SS7), GTPv1-C, RANAP}       & \lr{Diameter, GTPv2-C, S1AP} \\
\acro{TTI}           & ۲ میلی‌ثانیه (\acro{HSPA})، ۱۰ تا ۸۰ (\lr{R99}) & ۱ میلی‌ثانیه \\
\acro{IP} همیشه‌روشن & خیر (زمینهٔ \acro{PDP} بنا به تقاضا)  & بله (حامل پیش‌فرض هنگام اتصال) \\
بیشینهٔ فروسو (دستگاه رایج) & \lr{14--42\,Mbit/s}           & \lr{150--300\,Mbit/s} (نسخهٔ ۸) \\
\bottomrule
\end{tabularx}
\end{table}

\paragraph{مصالحه‌ها} طراحی تخت و تمام‌\acro{IP} شبکه را ساده کرد اما پیچیدگی را به جای دیگری برد: کیفیت
خدمات صدا اکنون باید با حامل‌ها و سیاست (\acro{PCRF}) مهندسی شود، \acro{eNB} کارکردها و حالت بیشتری
حمل می‌کند، تحرک بدون تحویل نرم به تحویل‌های سخت سریع متکی است و دورهٔ گذار به \acro{CSFB} و \acro{SRVCC}
نیاز داشت. جدایی صفحهٔ کنترل و کاربر که از پیش در \acro{EPC} وجود داشت (\acro{MME} در برابر دروازه‌ها) و با
\acro{CUPS} کامل شد، پیش‌درآمد هستهٔ خدمت‌محور \lr{5G} است.

% ---------------------------------------------------------------------
\section{نکته‌های کلیدی}
\begin{takeawaybox}
\begin{itemize}
  \item الزامات \acro{LTE}: بیشینهٔ ۱۰۰/۵۰ مگابیت بر ثانیه در ۲۰ مگاهرتز، تأخیر صفحهٔ کاربر کمتر از ۵
        میلی‌ثانیه، پهنای باند مقیاس‌پذیر، فقط بسته‌ای.
  \item \acro{EPS} = \acro{E-UTRAN} (فقط \acro{eNB}) + \acro{EPC} (\lr{MME, S-GW, P-GW, HSS, PCRF}).
  \item واسط‌های کنترلی از \lr{S1AP}، \lr{GTPv2-C} و \lr{Diameter} استفاده می‌کنند؛ دادهٔ کاربر از
        \lr{GTP-U} روی \lr{S1-U} و \lr{S5/S8}.
  \item \lr{S-GW} = لنگر تحرک؛ \lr{P-GW} = لنگر \acro{IP} و اعمال سیاست؛ \acro{MME} = فقط سیگنالینگ.
  \item حامل پیش‌فرض \acro{EPS} هنگام اتصال ساخته می‌شود: \acro{IP} همیشه‌روشن.
  \item بیشینهٔ نرخ $\approx\nu Q_m r_\mathrm{c} N_{\mathrm{RB}} N_{\mathrm{RE}}^{\mathrm{data}}/T_{\mathrm{sf}}$؛
        \acro{TTI} یک‌میلی‌ثانیه‌ای و حلقهٔ \acro{HARQ} هشت‌میلی‌ثانیه‌ای بر تأخیر رادیویی غالب‌اند.
  \item صدا: \acro{VoLTE} روی \acro{IMS} با حامل \lr{QCI 1}؛ \acro{CSFB} و \acro{SRVCC} برای هم‌کنشی با نسل دوم و سوم.
  \item \acro{CUPS} دروازه‌ها را به بخش‌های \lr{C} و \lr{U} که با \acro{PFCP} کنترل می‌شوند تقسیم می‌کند -- پلی به \lr{5G}.
\end{itemize}
\end{takeawaybox}

% ---------------------------------------------------------------------
\section{تمرین‌ها}

\begin{exercisebox}{1}{چرا دو دروازه؟}\label{exr:twogw}
توضیح دهید چرا \acro{EPC} هم \lr{S-GW} و هم \lr{P-GW} دارد. یک استقرار را مثال بزنید که در آن ترکیب
شده‌اند و یکی را که باید جدا باشند.
\end{exercisebox}

\begin{exercisebox}{1}{\lr{RNC} کجا رفت؟}\label{exr:rnc}
برای هر یک از کارکردهای زیر \acro{RNC} بگویید در \acro{LTE} کجا قرار دارد: پایان \acro{RRC}، رمزنگاری
دادهٔ کاربر، تصمیم تحویل، کنترل توان حلقهٔ بیرونی، لنگر تحرک به سوی هسته.
\end{exercisebox}

\begin{exercisebox}{2}{سربار پیشوند چرخه‌ای}\label{exr:cp}
در \acro{LTE} با پیشوند چرخه‌ای عادی، هر شیار ۰٫۵ میلی‌ثانیه‌ای ۷ نماد \acro{OFDM} به طول
$2048\,T_\mathrm{s}$ با $T_\mathrm{s} = 1/(15000\times2048)$\,\lr{s} دارد؛ نخستین پیشوند
$160\,T_\mathrm{s}$ و بقیه $144\,T_\mathrm{s}$ طول دارند. مدت پیشوندها را بر حسب میکروثانیه و کسری از زمان
را که صرف پیشوند می‌شود محاسبه کنید. بیشینهٔ گسترش تأخیر قابل تحمل بدون تداخل بین‌نمادی چقدر است؟
\end{exercisebox}

\begin{exercisebox}{2}{شبکهٔ انتقال \lr{VoLTE}}\label{exr:volte}
یک بستهٔ \acro{VoLTE} هر ۲۰ میلی‌ثانیه یک قاب گفتار ۳۳ بایتی را با ۱۲ بایت \acro{RTP}، ۸ بایت \acro{UDP} و
۲۰ بایت \lr{IPv4} حمل می‌کند. روی \lr{S1-U} در \lr{GTP-U/UDP/IPv4} بسته‌بندی می‌شود. نرخ بیت هر تماس روی
شبکهٔ انتقال و بازده انتقال \cref{eq:eff} را محاسبه کنید. چند تماس در یک پیوند \lr{100\,Mbit/s} با بار
۷۰٪ جا می‌گیرد؟
\end{exercisebox}

\begin{exercisebox}{3}{تأخیر فراسو}\label{exr:ullat}
\cref{eq:latency} را برای فراسوی \acro{UE}ای که مجوز ندارد گسترش دهید: باید منتظر فرصت بعدی درخواست
زمان‌بندی (\acro{SR}) با دورهٔ $T_{\mathrm{SR}}$ بماند، \acro{SR} را بفرستد، پس از $k_1$ میلی‌ثانیه مجوز
بگیرد و $k_2$ میلی‌ثانیه پس از مجوز ارسال کند. میانگین تأخیر را استخراج و برای $T_{\mathrm{SR}} = 10$\,\lr{ms}،
$k_1 = 3$\,\lr{ms}، $k_2 = 4$\,\lr{ms} و بقیهٔ مقادیر مانند \cref{sec:latency} ارزیابی کنید.
\end{exercisebox}

\begin{exercisebox}{3}{مسیریابی خانگی در برابر خروج محلی}\label{exr:roaming}
مشترکی از تهران در استانبول رومینگ است و از سروری در استانبول بازدید می‌کند. مسیر صفحهٔ کاربر و تأخیر
انتشار رفت‌وبرگشت را برای رومینگ با مسیریابی خانگی (\lr{S8} به \lr{P-GW} خانگی) و خروج محلی (\lr{P-GW}
میزبان) مقایسه کنید. فاصلهٔ دو شهر را ۲۰۰۰ کیلومتر، سرعت در فیبر را $2\times10^8$\,\lr{m/s} و ضریب مسیر را
۱٫۵ فرض کنید. هر گزینه برای صورت‌حساب و شنود قانونی چه معنایی دارد؟
\end{exercisebox}

\subsection*{پاسخ‌ها}
\addcontentsline{toc}{subsection}{پاسخ‌ها}

\begin{solutionbox}{\ref{exr:twogw}}
آن‌ها چیزهای متفاوتی را لنگر می‌کنند. \lr{S-GW} تحرک میان \acro{eNB}ها (و به نسل دوم/سوم) را لنگر و برای
\acro{UE}های بیکار بافر می‌کند؛ در هر لحظه یکی برای هر \acro{UE} وجود دارد و ممکن است با مکان تغییر کند.
\lr{P-GW} نشانی \acro{IP} و اتصال به یک \acro{PDN} را لنگر و سیاست را اعمال می‌کند؛ نباید در طول نشست تغییر
کند. ترکیب‌شده: یک اپراتور ملی کوچک، یک سایت، بدون رومینگ. جدا: رومینگ با مسیریابی خانگی (\lr{S-GW}
میزبان، \lr{P-GW} خانگی روی \lr{S8})، یا \acro{UE}ای با \acro{PDN} اینترنت روی یک \lr{P-GW} محلی و \acro{PDN}
\acro{IMS} روی یک \lr{P-GW} مرکزی.
\end{solutionbox}

\begin{solutionbox}{\ref{exr:rnc}}
پایان \acro{RRC}: \acro{eNB}. رمزنگاری دادهٔ کاربر: \acro{eNB} (\acro{PDCP}). تصمیم تحویل: \acro{eNB} مبدأ.
کنترل توان حلقهٔ بیرونی: به همان شکل لازم نیست -- \acro{LTE} از منابع زمان‌بندی‌شده و متعامد استفاده
می‌کند؛ کنترل توان فراسو (حلقهٔ باز به‌علاوهٔ اصلاحات حلقهٔ بسته) در \acro{eNB} است. لنگر تحرک به سوی هسته:
\lr{S-GW} (نقش لنگر \acro{Iu} در \acro{RNC}).
\end{solutionbox}

\begin{solutionbox}{\ref{exr:cp}}
$T_\mathrm{s} = 32.55$\,\lr{ns}. نخستین پیشوند: $160\,T_\mathrm{s} = \qty{5.21}{\micro\second}$؛ بقیه:
$144\,T_\mathrm{s} = \qty{4.69}{\micro\second}$؛ نماد مفید $2048\,T_\mathrm{s} = \qty{66.67}{\micro\second}$.
کسر پیشوند: $(160+6\times144)/15360 = 6.67\%$. گسترش تأخیر تا حدود $\qty{4.7}{\micro\second}$ (حدود ۱٫۴
کیلومتر اختلاف مسیر) جذب می‌شود؛ پیشوند گسترده ($512\,T_\mathrm{s} = \qty{16.7}{\micro\second}$) برای
سلول‌های بزرگ و پخش به کار می‌رود.
\end{solutionbox}

\begin{solutionbox}{\ref{exr:volte}}
بستهٔ \acro{IP}: $33+12+8+20 = 73$ بایت؛ روی \lr{S1-U}: $73+36 = 109$ بایت در هر ۲۰ میلی‌ثانیه
$= \qty{43.6}{\kilo\bit\per\second}$ (به‌علاوهٔ قاب‌بندی اترنت در عمل). بازده نسبت به محمولهٔ گفتار:
$33/109 = 30\%$؛ نسبت به بستهٔ \acro{IP} کاربر طبق \cref{eq:eff}: $73/109 = 67\%$. تعداد تماس:
$0.7\times100\times10^6/43\,600 \approx 1600$ در هر جهت.
\end{solutionbox}

\begin{solutionbox}{\ref{exr:ullat}}
میانگین انتظار برای فرصت \acro{SR} برابر $T_{\mathrm{SR}}/2$ است، سپس ارسال \acro{SR} (۱ میلی‌ثانیه)، مجوز پس
از $k_1$، داده $k_2$ پس از مجوز و سپس جمله‌های مشابه فروسو:
$D_{\mathrm{UL}} = T_{\mathrm{SR}}/2 + T_{\mathrm{TTI}} + k_1 + k_2 + T_{\mathrm{TTI}} + D_{\mathrm{eNB}} +
p_{\mathrm{re}}T_{\mathrm{HARQ}} + D_{\mathrm{tr}} = 5 + 1 + 3 + 4 + 1 + 1 + 0.8 + 0.5 = \qty{16.3}{\milli\second}$.
چرخهٔ درخواست زمان‌بندی غالب است؛ پیش‌زمان‌بندی یا دوره‌های کوتاه \acro{SR} آن را کاهش می‌دهند و \lr{5G}
مجوزهای پیکربندی‌شده‌ای می‌افزاید که آن را حذف می‌کنند.
\end{solutionbox}

\begin{solutionbox}{\ref{exr:roaming}}
مسیریابی خانگی: \acro{UE} در استانبول $\to$ \lr{S-GW} میزبان $\to$ \lr{S8} به \lr{P-GW} تهران $\to$ بازگشت
به سرور استانبول: دو بار پیمودن $1.5\times2000 = 3000$ کیلومتر در هر جهت، یعنی ۱۲۰۰۰ کیلومتر در هر
رفت‌وبرگشت یا ۶۰ میلی‌ثانیه انتشار. خروج محلی: عملاً محلی، بسیار کمتر از ۱ میلی‌ثانیه. مسیریابی خانگی
صورت‌حساب، سیاست و شنود قانونی را در شبکهٔ خانگی نگه می‌دارد (برای اپراتور خانگی ساده است)؛ خروج محلی
مستلزم آن است که شبکهٔ میزبان سیاست اپراتور خانگی را اعمال کند (از طریق \lr{S9} یا توافق‌های رومینگ)
و شنود را به‌صورت محلی انجام دهد.
\end{solutionbox}

% ---------------------------------------------------------------------
\section{مطالعهٔ بیشتر}
توصیف کلی \acro{E-UTRAN} در \cite{ts36300} و رویه‌های \acro{EPC} در \cite{ts23401} آمده‌اند؛ الزامات
\acro{LTE} در \cite{tr25913} است. کنترل سیاست و صورت‌حساب در \cite{ts23203}، \acro{CUPS} در \cite{ts23214}
همراه با \acro{PFCP} در \cite{ts29244}، \acro{CSFB} در \cite{ts23272} و \acro{SRVCC} در \cite{ts23216}
مشخص شده‌اند. \lr{Sesia} و همکاران \cite{sesia2011lte} و \lr{Olsson} و همکاران \cite{olsson2013epc} به ترتیب
کتاب‌های پیشنهادی برای سمت رادیو و هسته هستند.

\begin{latin}
\printbibliography[heading=subbibliography,title={References}]
\end{latin}
\end{document}
